% =====================================================================
% Premessa: come leggere il manuale (legenda dei riquadri e fonti)
% =====================================================================
\chapter*{Come leggere questo manuale}
\addcontentsline{toc}{chapter}{Come leggere questo manuale}

Questo manuale è scritto da studenti, per studenti: \textbf{non è materiale ufficiale del corso} e non
sostituisce le slide, i notebook e le indicazioni dei docenti. Nasce dalle registrazioni delle
lezioni (trascritte automaticamente e poi revisionate), integrate con slide e materiale dei laboratori.

\section*{Le parti}
\begin{description}[leftmargin=2em,style=nextline]
  \item[Diario delle lezioni] una voce per lezione, in ordine cronologico: sintesi, avvisi pratici,
    contenuti in breve e punti segnalati per l'esame. Serve a seguire l'avanzamento del corso.
  \item[Manuale del corso] un capitolo per argomento (dalle basi del deep learning alle architetture,
    agli embedding e all'anomaly detection), che unisce lezioni, slide e notebook in un'esposizione
    ordinata. Rimanda alle lezioni del diario.
  \item[Laboratori] come si scrive il report e, per ogni lab, contesto, consegna, metodo e insidie.
    \textbf{Non contiene soluzioni}: i lab sono valutati e vanno svolti dal proprio gruppo.
  \item[Schemi di ripasso] mappe concettuali e tabelle ``quale scelta quando''.
  \item[Quiz ed esame] formato dell'esame, quiz a scelta multipla e vero/falso, domande aperte ed
    esercizi su pipeline e codice, con soluzioni e rimando alla lezione.
  \item[Appendice] glossario e acronimi.
\end{description}

\section*{Legenda dei riquadri}

\begin{esame}
Punto che il docente ha indicato come rilevante per l'esame.
\end{esame}

\begin{integrazione}
Contenuto aggiunto dai redattori: \textbf{non detto a lezione e non presente nelle slide}
(spiegazioni, esempi, riferimenti bibliografici in più). Utile, ma da trattare con cautela all'esame.
\end{integrazione}

\begin{verifica}
Informazione non confermata: passaggio poco chiaro nella registrazione, dato ricostruito o calcolato.
\end{verifica}

\begin{nota}
Nota generica (fonti usate, avvertenze, collegamenti).
\end{nota}

\begin{soluzione}
Soluzione di quiz, simulazioni e laboratori.
\end{soluzione}

Riquadri didattici, con contenuti presi dal corso:

\begin{definizione}[termine]
Definizione di un concetto.
\end{definizione}

\begin{esempio}
Esempio fatto a lezione o nelle slide.
\end{esempio}

\begin{attenzione}
Errore comune o insidia segnalata dal docente.
\end{attenzione}

\begin{sintesi}
Riepilogo dei punti chiave.
\end{sintesi}

\dacompletare{Parte ancora da scrivere (per esempio un argomento non ancora trattato a lezione).}

I termini tecnici restano in inglese, come nelle slide e all'esame, con la traduzione alla prima
occorrenza; il glossario in appendice raccoglie le corrispondenze.

\section*{Fonti}
Registrazioni delle lezioni (Prof.~Danilo Giordano; Prof.~Matteo Boffa per NLP e ultima parte del
corso), slide, notebook e testi dei laboratori pubblicati sul portale della didattica. Slide, notebook,
registrazioni e dataset non sono inclusi nel repository.
